\documentclass[12pt,a4paper]{article}
\usepackage[utf8]{inputenc}
\usepackage[swedish]{babel}
\usepackage{geometry}
\usepackage{graphicx}
\usepackage{xcolor}
\usepackage{hyperref}
\usepackage{listings}
\usepackage{tikz}
\usepackage{enumitem}
\usepackage{booktabs}
\usepackage{array}
\usepackage{verbatim}

\geometry{margin=2.5cm}

% Define brand colors
\definecolor{brand}{RGB}{16,185,129}
\definecolor{brandlight}{RGB}{209,250,229}
\definecolor{branddark}{RGB}{4,120,87}

\title{\textbf{Administratörspanel}\\
\large Teknisk Specifikation och UI/UX Design}
\author{Roaring Analytics AB\\
Utvecklad för LIA2-praktik och framtida produktionsmiljö}
\date{Version 1.0 -- \today}

\begin{document}

\maketitle
\tableofcontents
\newpage

% ═══════════════════════════════════════════════════════════════════════
\section{Översikt}
% ═══════════════════════════════════════════════════════════════════════

\subsection{Syfte}
Administratörspanelen är en central plattform för Roaring Analytics AB:s personal att övervaka, hantera och analysera alla användare i systemet. Panelen ger insikt i:

\begin{itemize}[leftmargin=2cm]
  \item \textbf{Fraud Detection:} Aggregerad statistik över upptäckta bedrägerier för rapportering till Finanspolisen
  \item \textbf{Support:} Hantera support-ärenden från revisionsbyråer
  \item \textbf{Fakturering:} Månadsvis fakturering och betalningsstatus
  \item \textbf{Email:} Kommunikation med användare (enskilt och massutskick)
  \item \textbf{Användare:} Översikt över alla registrerade revisionsbyråer
\end{itemize}

\subsection{Målgrupp}
\begin{itemize}
  \item \textbf{Administratörer:} Anställda på Roaring Analytics AB
  \item \textbf{Säljteam:} Använder fraud detection-rapporter för demos till Finanspolisen
  \item \textbf{Support:} Customer success team som hjälper revisionsbyråer
\end{itemize}

\subsection{Teknisk Stack}
\begin{itemize}
  \item \textbf{Frontend:} React 18 + Tailwind CSS
  \item \textbf{Backend:} FastAPI (Python) + PostgreSQL
  \item \textbf{Email:} SendGrid API (production) / Mock (dev)
  \item \textbf{Routing:} React Router v6 (\texttt{/admin})
  \item \textbf{Framtida uppgradering:} Themesberg React/Tailwind Admin Template (\$89)
\end{itemize}

\subsection{Separation från Användar-UI}
Administratörspanelen är \textbf{helt separerad} från byråchefernas onboarding-gränssnitt:

\begin{table}[h]
\centering
\begin{tabular}{|l|l|}
\hline
\textbf{Byråchef (Användare)} & \textbf{Administrator (Roaring)} \\
\hline
Route: \texttt{/dokumentation, /underlag, /bokforing} & Route: \texttt{/admin} \\
Sidebar: Onboarding-steg & Sidebar: Admin-funktioner \\
Färgschema: Grön (brand) & Färgschema: Lila (admin) \\
Behörighet: JWT token för egen byrå & Behörighet: Admin-roll i JWT \\
\hline
\end{tabular}
\end{table}

\vspace{1cm}
\textbf{OBS:} I framtida version kommer AdminDashboard att ha sin egen layout med separat sidebar och header.

% ═══════════════════════════════════════════════════════════════════════
\section{Flik 1: Översikt (Overview)}
% ═══════════════════════════════════════════════════════════════════════

\subsection{Syfte}
Ge en snabb överblick av systemets status med key metrics och senaste aktiviteter.

\subsection{Komponenter}

\subsubsection{Statistikkort (Stats Cards)}
4 kort med centrala KPI:er:

\begin{enumerate}
  \item \textbf{Totalt Användare}
  \begin{itemize}
    \item Antal: 450
    \item Ikon:  (brand-grön bakgrund)
    \item Källa: \texttt{SELECT COUNT(*) FROM users}
  \end{itemize}
  
  \item \textbf{Aktiva Onboardings}
  \begin{itemize}
    \item Antal: 23
    \item Ikon: \texttt{Clock icon} (blå bakgrund)
    \item Definition: Användare som startat men ej avslutat onboarding
    \item Källa: \texttt{SELECT COUNT(*) FROM onboarding\_sessions WHERE status='active'}
  \end{itemize}
  
  \item \textbf{Flaggade Ärenden}
  \begin{itemize}
    \item Antal: 127
    \item Ikon: \texttt{Alert icon} (röd bakgrund)
    \item Definition: Fraud alerts med risk score > 40
    \item Källa: \texttt{SELECT COUNT(*) FROM fraud\_alerts WHERE created\_at > NOW() - INTERVAL '30 days'}
  \end{itemize}
  
  \item \textbf{Nya Ansökningar}
  \begin{itemize}
    \item Antal: 8
    \item Ikon: \texttt{Check icon} (grön bakgrund)
    \item Definition: Nya registreringar senaste 7 dagarna
    \item Källa: \texttt{SELECT COUNT(*) FROM users WHERE created\_at > NOW() - INTERVAL '7 days'}
  \end{itemize}
\end{enumerate}

\subsubsection{Senaste Fraud Alerts}
Panel med top 3 senaste fraud alerts:

\begin{itemize}
  \item \textbf{Layout:} Kort med risk score badge (cirkel), företagsnamn, beskrivning, datum, status
  \item \textbf{Färgkodning:} 
  \begin{itemize}
    \item Risk > 80: Röd bakgrund
    \item Risk 60-80: Orange bakgrund
    \item Risk 40-60: Gul bakgrund
  \end{itemize}
  \item \textbf{Status badges:} Ny (blå), Under granskning (gul), Avslutad (grön)
  \item \textbf{Interaktion:} Klick på "Visa alla fraud alerts →" navigerar till Fraud Detection-fliken
\end{itemize}

\subsubsection{Senaste Support-ärenden}
Panel med top 3 senaste support-ärenden:

\begin{itemize}
  \item \textbf{Layout:} Kort med ämne, företag, kontaktperson, prioritet, status
  \item \textbf{Prioritet badges:} Hög (röd), Medium (gul), Låg (grön)
  \item \textbf{Status badges:} Ny (blå), Under behandling (gul), Avslutad (grön)
  \item \textbf{Interaktion:} Klick på "Visa alla support-ärenden →" navigerar till Support-fliken
\end{itemize}

% ═══════════════════════════════════════════════════════════════════════
\section{Flik 2: Fraud Detection}
% ═══════════════════════════════════════════════════════════════════════

\subsection{Syfte}
Central plats för att övervaka fraud detection-resultat och generera rapporter för Finanspolisen.

\subsection{Key Metrics}

\begin{table}[h]
\centering
\begin{tabular}{|l|r|p{6cm}|}
\hline
\textbf{Metric} & \textbf{Värde} & \textbf{Definition} \\
\hline
Total Alerts (30 dagar) & 127 & Alla fraud alerts senaste månaden \\
Hög Risk (>80) & 23 & Alerts med risk score över 80 \\
Detektionsgrad & 28.2\% & Andel användare med minst 1 alert \\
\hline
\end{tabular}
\end{table}

\subsection{Detektionskategorier}
Stapeldiagram med 5 kategorier:

\begin{enumerate}
  \item \textbf{Privatkonsumtion:} 57 alerts (45\%)
  \begin{itemize}
    \item \textit{Exempel:} ICA, H\&M, Willys bokförda på företaget
    \item \textit{Färg:} Orange (\texttt{bg-orange-500})
  \end{itemize}
  
  \item \textbf{Verksamhetskongruens:} 38 alerts (30\%)
  \begin{itemize}
    \item \textit{Exempel:} Cykelförsäljare köper fiskeredskap
    \item \textit{Färg:} Röd (\texttt{bg-red-500})
  \end{itemize}
  
  \item \textbf{Cirkulära betalningar:} 19 alerts (15\%)
  \begin{itemize}
    \item \textit{Exempel:} Betalningar mellan närstående företag
    \item \textit{Färg:} Gul (\texttt{bg-yellow-500})
  \end{itemize}
  
  \item \textbf{Konkurskontroll:} 8 alerts (6\%)
  \begin{itemize}
    \item \textit{Exempel:} Betalning till avregistrerat företag
    \item \textit{Färg:} Blå (\texttt{bg-blue-500})
  \end{itemize}
  
  \item \textbf{Leveransadress:} 5 alerts (4\%)
  \begin{itemize}
    \item \textit{Exempel:} Leveransadress matchar ej företagsadress
    \item \textit{Färg:} Lila (\texttt{bg-purple-500})
  \end{itemize}
\end{enumerate}

\subsection{Exportfunktion: Rapport för Finanspolisen}
Stor grön knapp: \textbf{"Generera Rapport för Finanspolisen (PDF)"}

\subsubsection{Rapportinnehåll}
\begin{itemize}
  \item \textbf{Aggregerad statistik:} Total antal alerts, detektionsgrad, fördelning per kategori
  \item \textbf{Anonymiserade case studies:} 
  \begin{itemize}
    \item Case A: Cykelförsäljare - fiskeredskap (85/100)
    \item Case B: Privata ICA-inköp (78/100)
    \item Case C: Cirkulära betalningar (92/100)
  \end{itemize}
  \item \textbf{Trendanalys:} Utveckling över tid (senaste 12 månader)
  \item \textbf{ROI för revisionsbyråer:} Sparad tid per månad tack vare automatisk detektion
  \item \textbf{False positive rate:} Andel alerts som markerats som "falskt alarm"
\end{itemize}

\subsubsection{Anonymiseringsregler}
\textbf{KRITISKT:} Rapporten får ALDRIG innehålla:
\begin{itemize}
  \item Verkliga företagsnamn
  \item Organisationsnummer (endast maskerade: 559XXX-XXXX)
  \item Personnummer
  \item Exakta adresser
\end{itemize}

Använd istället:
\begin{itemize}
  \item "Anonymiserat företag A, B, C..."
  \item "Case 001, 002, 003..."
  \item "Verksamhetstyp: Cykelförsäljare" (ej företagsnamn)
\end{itemize}

\subsection{Tabell: Alla Fraud Alerts}
Full tabell med kolumner:
\begin{itemize}
  \item Företag (anonymiserat)
  \item Org.nr (maskerat)
  \item Kategori
  \item Beskrivning
  \item Risk Score (färgkodad cirkel)
  \item Datum
  \item Status (badge)
\end{itemize}

% ═══════════════════════════════════════════════════════════════════════
\section{Flik 3: Support}
% ═══════════════════════════════════════════════════════════════════════

\subsection{Syfte}
Hantera support-ärenden från revisionsbyråer som behöver hjälp med systemet.

\subsection{Ärendekö}
Lista på alla support-ärenden med följande information:

\begin{itemize}
  \item \textbf{Ämne:} Kort beskrivning av problemet
  \item \textbf{Meddelande:} Fullständig text från användaren
  \item \textbf{Företag:} Revisionsbyrån som skickat ärendet
  \item \textbf{Kontaktperson:} Namn på byråchef
  \item \textbf{Email:} Kontaktuppgifter
  \item \textbf{Prioritet:} Hög / Medium / Låg
  \item \textbf{Status:} Ny / Under behandling / Avslutad
  \item \textbf{Skapad:} Datum och tid
\end{itemize}

\subsection{Interaktioner}
\begin{enumerate}
  \item \textbf{Svara-knapp:} Öppnar Email-fliken med förifylld mottagare
  \item \textbf{Filtrera:} Efter status och prioritet
  \item \textbf{Sortera:} Efter datum (nyast först som standard)
  \item \textbf{Markera som behandlad:} Ändrar status till "Under behandling" eller "Avslutad"
\end{enumerate}

\subsection{Mock Data (Demo)}
3 exempel-ärenden:

\begin{enumerate}
  \item \textbf{Revision Stockholm AB}
  \begin{itemize}
    \item Kontakt: Anna Andersson
    \item Ämne: "Fråga om fraud alert - företag A"
    \item Prioritet: Hög
    \item Status: Ny
  \end{itemize}
  
  \item \textbf{Ekonomibyrån Väst}
  \begin{itemize}
    \item Kontakt: Erik Eriksson
    \item Ämne: "Problem med uppladdning av SIE-filer"
    \item Prioritet: Medium
    \item Status: Under behandling
  \end{itemize}
  
  \item \textbf{Nordisk Revision}
  \begin{itemize}
    \item Kontakt: Maria Nilsson
    \item Ämne: "Fakturafråga"
    \item Prioritet: Låg
    \item Status: Avslutad
  \end{itemize}
\end{enumerate}

% ═══════════════════════════════════════════════════════════════════════
\section{Flik 4: Fakturering}
% ═══════════════════════════════════════════════════════════════════════

\subsection{Syfte}
Översikt över månadsvis fakturering till revisionsbyråer baserat på antal användare.

\subsection{Prismodell}
\begin{itemize}
  \item \textbf{Bas:} 950 kr/månad per användare
  \item \textbf{Volymrabatt:}
  \begin{itemize}
    \item 3-5 användare: -10\% (855 kr/användare)
    \item 6-10 användare: -20\% (760 kr/användare)
    \item 11+ användare: -30\% (665 kr/användare)
  \end{itemize}
  \item \textbf{Fakturering:} Månadsvis i efterskott
  \item \textbf{Betalningsvillkor:} 30 dagar netto
\end{itemize}

\subsection{Fakturatabell}
Kolumner:
\begin{itemize}
  \item \textbf{Faktura-ID:} INV-2025-001, INV-2025-002, etc.
  \item \textbf{Företag:} Revisionsbyrån
  \item \textbf{Period:} 2025-10 (oktober 2025)
  \item \textbf{Användare:} Antal aktiva användare under perioden
  \item \textbf{Belopp:} Total summa i SEK
  \item \textbf{Förfallodatum:} Datum
  \item \textbf{Status:} Betald (grön) / Obetald (röd)
  \item \textbf{Betald datum:} Om betald
\end{itemize}

\subsection{Mock Data (Demo)}
\begin{table}[h]
\centering
\small
\begin{tabular}{|l|l|c|c|r|l|}
\hline
\textbf{ID} & \textbf{Företag} & \textbf{Användare} & \textbf{Belopp} & \textbf{Status} & \textbf{Förfallodatum} \\
\hline
INV-2025-001 & Revision Stockholm AB & 3 & 2 950 kr & Betald & 2025-11-15 \\
INV-2025-002 & Ekonomibyrån Väst & 2 & 1 950 kr & Betald & 2025-11-15 \\
INV-2025-003 & Nordisk Revision & 5 & 4 950 kr & Obetald & 2025-11-15 \\
\hline
\end{tabular}
\end{table}

\subsection{Exportfunktion}
Knapp: \textbf{"Exportera till Excel"}
\begin{itemize}
  \item Genererar \texttt{.xlsx}-fil med alla fakturor
  \item Inkluderar summeringsrad
  \item Filnamn: \texttt{Fakturering\_2025-10.xlsx}
\end{itemize}

\subsection{Total Summa}
Nederst i tabellen:
\begin{itemize}
  \item \textbf{Totalt belopp:} Summa av alla fakturor (9 850 kr i mock data)
  \item \textbf{Betald:} 4 900 kr
  \item \textbf{Obetald:} 4 950 kr
\end{itemize}

% ═══════════════════════════════════════════════════════════════════════
\section{Flik 5: Email}
% ═══════════════════════════════════════════════════════════════════════

\subsection{Syfte}
Skicka email till användare, antingen enskilt eller som massutskick till alla 450 användare.

\subsection{Användningsområden}
\begin{itemize}
  \item \textbf{Support-svar:} Svara på support-ärenden
  \item \textbf{Notifikationer:} Påminnelser om ofullständiga onboardings
  \item \textbf{Fraud alerts:} Automatiska varningar vid hög risk score
  \item \textbf{Massutskick:} Nyhetsbrev, produktuppdateringar, systemunderhåll
  \item \textbf{REMOTE\_ONBOARDING:} Skicka länk till extern uppladdning för stora datafiler
\end{itemize}

\subsection{Formulär}

\subsubsection{Email-typ (Radio buttons)}
\begin{itemize}
  \item \textbf{Enskild användare:} Skicka till en specifik mottagare
  \item \textbf{Massutskick:} Skicka till alla 450 användare
\end{itemize}

\subsubsection{Fält}
\begin{enumerate}
  \item \textbf{Till:} (endast för enskild användare)
  \begin{itemize}
    \item Input: Email-adress
    \item Validering: Måste vara giltig email
    \item Placeholder: "mottagare@example.com"
  \end{itemize}
  
  \item \textbf{Ämne:}
  \begin{itemize}
    \item Input: Textsträng
    \item Max 100 tecken
    \item Obligatoriskt
  \end{itemize}
  
  \item \textbf{Meddelande:}
  \begin{itemize}
    \item Textarea: 8 rader
    \item Stöd för line breaks
    \item Framtida version: Rich text editor (från Themesberg-template)
  \end{itemize}
\end{enumerate}

\subsubsection{Knappar}
\begin{itemize}
  \item \textbf{Skicka Email:} Grön knapp, fullbredds
  \item \textbf{Rensa:} Grå knapp, rensa formulär
\end{itemize}

\subsection{Integration: SendGrid API}

\subsubsection{Production Environment}
\begin{lstlisting}[basicstyle=\small\ttfamily]
import sendgrid
from sendgrid.helpers.mail import Mail

def send_email(to_email, subject, message):
    sg = sendgrid.SendGridAPIClient(
        api_key=os.environ.get('SENDGRID_API_KEY')
    )
    
    mail = Mail(
        from_email='admin@roaringanalytics.se',
        to_emails=to_email,
        subject=subject,
        html_content=message
    )
    
    response = sg.send(mail)
    return response.status_code
\end{lstlisting}

\subsubsection{Development Environment}
Mock-funktionalitet:
\begin{itemize}
  \item \texttt{alert()} med sammanfattning av email
  \item Loggas till console
  \item Ingen faktisk email skickas
\end{itemize}

\subsection{Säkerhet}
\begin{itemize}
  \item \textbf{Rate limiting:} Max 100 email/timme per admin
  \item \textbf{Massutskick:} Kräver bekräftelse-dialog innan skickning
  \item \textbf{HTML sanitization:} Alla användare-input saniteras för XSS
  \item \textbf{Email validering:} Regex-check för giltiga email-adresser
\end{itemize}

% ═══════════════════════════════════════════════════════════════════════
\section{Teknisk Implementation}
% ═══════════════════════════════════════════════════════════════════════

\subsection{Filstruktur}
\begin{verbatim}
src/
  components/
    Admin/
      AdminDashboard.jsx (huvudkomponent, 870+ rader)
      AdminSidebar.jsx (framtida - separat sidebar)
      AdminHeader.jsx (framtida - separat header)
      components/
        StatsCard.jsx
        FraudAlertCard.jsx
        SupportTicketCard.jsx
        InvoiceTable.jsx
        EmailForm.jsx
  services/
    adminApi.js (API-anrop för admin-funktioner)
    emailService.js (SendGrid integration)
\end{verbatim}

\subsection{State Management}
\begin{lstlisting}[basicstyle=\small\ttfamily]
const [activeTab, setActiveTab] = useState('overview');
// 'overview', 'fraud', 'support', 'invoices', 'email'

const [emailForm, setEmailForm] = useState({
  to: '',
  subject: '',
  message: '',
  type: 'single' // 'single' or 'mass'
});
\end{lstlisting}

\subsection{API Endpoints}

\subsubsection{Översikt}
\begin{itemize}
  \item \texttt{GET /api/admin/stats} - Hämta KPI-statistik
  \item \texttt{GET /api/admin/recent-alerts} - Senaste fraud alerts (top 3)
  \item \texttt{GET /api/admin/recent-tickets} - Senaste support-ärenden (top 3)
\end{itemize}

\subsubsection{Fraud Detection}
\begin{itemize}
  \item \texttt{GET /api/admin/fraud/overview} - Aggregerad statistik
  \item \texttt{GET /api/admin/fraud/categories} - Kategorifördelning
  \item \texttt{GET /api/admin/fraud/alerts} - Alla alerts med filtrering
  \item \texttt{POST /api/admin/fraud/export-report} - Generera PDF-rapport
\end{itemize}

\subsubsection{Support}
\begin{itemize}
  \item \texttt{GET /api/admin/support/tickets} - Alla support-ärenden
  \item \texttt{PATCH /api/admin/support/tickets/:id} - Uppdatera status
  \item \texttt{POST /api/admin/support/reply} - Skicka svar via email
\end{itemize}

\subsubsection{Fakturering}
\begin{itemize}
  \item \texttt{GET /api/admin/invoices} - Alla fakturor
  \item \texttt{GET /api/admin/invoices/export} - Exportera till Excel
  \item \texttt{PATCH /api/admin/invoices/:id/status} - Markera som betald
\end{itemize}

\subsubsection{Email}
\begin{itemize}
  \item \texttt{POST /api/admin/email/send} - Skicka enskild email
  \item \texttt{POST /api/admin/email/mass-send} - Massutskick till alla användare
\end{itemize}

\subsection{Autentisering \& Behörighet}
\begin{itemize}
  \item \textbf{JWT Token:} Måste innehålla \texttt{role: 'admin'}
  \item \textbf{Middleware:} \texttt{@require\_admin} decorator på alla admin-endpoints
  \item \textbf{Frontend:} Dölj admin-ikon i header för icke-admins
  \item \textbf{Route Guard:} Redirect till \texttt{/unauthorized} om användare försöker nå \texttt{/admin}
\end{itemize}

% ═══════════════════════════════════════════════════════════════════════
\section{Framtida Uppgraderingar}
% ═══════════════════════════════════════════════════════════════════════

\subsection{Themesberg React/Tailwind Template}
Efter köp av template (25:e november 2025):

\subsubsection{Fördelar}
\begin{itemize}
  \item \textbf{Gmail-liknande email-interface:}
  \begin{itemize}
    \item Inkorgen med trådar
    \item Rich text editor (WYSIWYG)
    \item Bifogade filer
    \item Mappar och etiketter
  \end{itemize}
  
  \item \textbf{Professionella komponenter:}
  \begin{itemize}
    \item Chart.js/Recharts integration för avancerade grafer
    \item Drag-and-drop kanban board för support-ärenden
    \item Advanced data tables med sortering och filtrering
    \item Dark mode support
  \end{itemize}
  
  \item \textbf{Separat layout:}
  \begin{itemize}
    \item AdminSidebar med collapsible meny
    \item AdminHeader med notifications bell
    \item Breadcrumbs navigation
    \item User profile dropdown
  \end{itemize}
\end{itemize}

\subsubsection{Migration Plan}
\begin{enumerate}
  \item Installera Themesberg-template i \texttt{src/templates/themesberg/}
  \item Skapa \texttt{AdminLayout.jsx} som wrapper för admin-sidor
  \item Migrera existerande komponenter till nya layout
  \item Integrera Themesberg's EmailInbox-komponent
  \item Anpassa färgschema till Roaring's brand colors
  \item Testa responsiv design på mobil/tablet
\end{enumerate}

\subsection{Ytterligare Features}

\subsubsection{Real-time Notifications}
\begin{itemize}
  \item WebSocket-anslutning för live updates
  \item Toast notifications för nya support-ärenden
  \item Desktop notifications (browser API)
\end{itemize}

\subsubsection{Användare-hantering}
\begin{itemize}
  \item Lista alla användare med sök och filtrering
  \item Se användares onboarding-progress
  \item Impersonate-funktion för felsökning
  \item Skapa ny användare manuellt
\end{itemize}

\subsubsection{Avancerad Rapportering}
\begin{itemize}
  \item Custom datum-intervall för fraud detection
  \item Exportera till Excel/CSV
  \item Schemalagda rapporter (veckovis/månadsvis email)
  \item Jämförelse mellan olika perioder
\end{itemize}

\subsubsection{Aktivitetslogg}
\begin{itemize}
  \item Audit trail för alla admin-åtgärder
  \item Vem skickade vilken email och när
  \item Ändringar i support-ärenden
  \item Faktura-uppdateringar
\end{itemize}

% ═══════════════════════════════════════════════════════════════════════
\section{Mock Data för LIA-Demo}
% ═══════════════════════════════════════════════════════════════════════

\subsection{Fraud Alerts (Anonymiserade)}
\begin{table}[h]
\centering
\small
\begin{tabular}{|l|l|c|l|}
\hline
\textbf{Företag} & \textbf{Kategori} & \textbf{Risk} & \textbf{Beskrivning} \\
\hline
Anonymiserat A & Verksamhetskongruens & 85 & Cykelförsäljare - fiskeredskap \\
Anonymiserat B & Privatkonsumtion & 78 & ICA-inköp bokförda \\
Anonymiserat C & Cirkulära betalningar & 92 & Närstående företag \\
Anonymiserat D & Konkurskontroll & 65 & Avregistrerat leverantör \\
\hline
\end{tabular}
\end{table}

\subsection{Support-ärenden}
\begin{enumerate}
  \item \textbf{Revision Stockholm AB} - Anna Andersson
  \begin{itemize}
    \item Ämne: Fråga om fraud alert
    \item Prioritet: Hög
    \item Status: Ny
  \end{itemize}
  
  \item \textbf{Ekonomibyrån Väst} - Erik Eriksson
  \begin{itemize}
    \item Ämne: Problem med SIE-uppladdning
    \item Prioritet: Medium
    \item Status: Under behandling
  \end{itemize}
  
  \item \textbf{Nordisk Revision} - Maria Nilsson
  \begin{itemize}
    \item Ämne: Fakturafråga
    \item Prioritet: Låg
    \item Status: Avslutad
  \end{itemize}
\end{enumerate}

\subsection{Fakturor}
\begin{itemize}
  \item INV-2025-001: Revision Stockholm AB, 2 950 kr, Betald
  \item INV-2025-002: Ekonomibyrån Väst, 1 950 kr, Betald
  \item INV-2025-003: Nordisk Revision, 4 950 kr, Obetald
\end{itemize}

% ═══════════════════════════════════════════════════════════════════════
\section{Säkerhet \& GDPR}
% ═══════════════════════════════════════════════════════════════════════

\subsection{Dataskydd}
\begin{itemize}
  \item \textbf{Personuppgifter:} Admin-panelen innehåller känslig data
  \item \textbf{Åtkomstloggning:} Alla admin-åtgärder loggas med timestamp och user\_id
  \item \textbf{Kryptering:} All data krypteras i transit (HTTPS) och i vila (PostgreSQL encryption)
  \item \textbf{Datalagring:} Fraud alerts sparas i 7 år (enligt Bokföringslagen)
\end{itemize}

\subsection{Anonymisering}
\textbf{KRITISKT:} Vid export av rapporter till Finanspolisen:
\begin{itemize}
  \item Maskera alla företagsnamn: "Anonymiserat företag A, B, C..."
  \item Maskera org.nr: 559XXX-XXXX (visa endast första 3 siffror)
  \item Ta bort personnummer helt
  \item Generalisera adresser: "Stockholm" istället för "Bangårdsgatan 26, Östersund"
\end{itemize}

\subsection{Rollbaserad Åtkomst}
\begin{table}[h]
\centering
\begin{tabular}{|l|c|c|c|}
\hline
\textbf{Funktion} & \textbf{Admin} & \textbf{Support} & \textbf{Byråchef} \\
\hline
Översikt & \texttt{Check icon} & \texttt{Check icon} & Nej \\
Fraud Detection & \texttt{Check icon} & \texttt{Check icon} (read-only) & Nej \\
Support-ärenden & \texttt{Check icon} & \texttt{Check icon} & Nej \\
Fakturering & \texttt{Check icon} & Nej & Nej \\
Email (massutskick) & \texttt{Check icon} & Nej & Nej \\
Email (enskild) & \texttt{Check icon} & \texttt{Check icon} & Nej \\
\hline
\end{tabular}
\end{table}

% ═══════════════════════════════════════════════════════════════════════
\section{Testplan}
% ═══════════════════════════════════════════════════════════════════════

\subsection{Manuell Testning}
\begin{enumerate}
  \item \textbf{Navigering:}
  \begin{itemize}
    \item Klicka på lila admin-ikon i header → ska navigera till \texttt{/admin}
    \item Testa alla 5 flikar (Översikt, Fraud Detection, Support, Fakturering, Email)
    \item Verifiera att aktiv flik highlightas med brand-grön färg
  \end{itemize}
  
  \item \textbf{Översikt:}
  \begin{itemize}
    \item Verifiera att alla 4 statistikkort visar rätt värden
    \item Klicka "Visa alla fraud alerts" → ska navigera till Fraud Detection-fliken
    \item Klicka "Visa alla support-ärenden" → ska navigera till Support-fliken
  \end{itemize}
  
  \item \textbf{Fraud Detection:}
  \begin{itemize}
    \item Verifiera stapeldiagram med 5 kategorier
    \item Klicka "Generera Rapport för Finanspolisen" → ska visa alert med rapportinnehåll
    \item Kontrollera att alla företag är anonymiserade i tabellen
  \end{itemize}
  
  \item \textbf{Support:}
  \begin{itemize}
    \item Verifiera att alla 3 support-ärenden visas
    \item Klicka "Svara"-knapp → ska öppna Email-fliken med förifylld mottagare
    \item Verifiera färgkodning av prioritet och status badges
  \end{itemize}
  
  \item \textbf{Fakturering:}
  \begin{itemize}
    \item Kontrollera att total summa är korrekt (9 850 kr)
    \item Verifiera status badges (grön för betald, röd för obetald)
    \item Klicka "Exportera till Excel" → ska visa alert (mock)
  \end{itemize}
  
  \item \textbf{Email:}
  \begin{itemize}
    \item Testa båda email-typer (enskild och massutskick)
    \item Fyll i formulär och skicka → ska visa alert med sammanfattning
    \item Klicka "Rensa" → ska nollställa formulär
    \item Verifiera att "Till"-fält döljs för massutskick
  \end{itemize}
\end{enumerate}

\subsection{Responsiv Design}
\begin{itemize}
  \item \textbf{Desktop (>1280px):} Full bredd, alla kolumner synliga
  \item \textbf{Tablet (768-1280px):} 2 kolumner för statistikkort, scroll för tabeller
  \item \textbf{Mobil (<768px):} 1 kolumn, tabs ska vara scrollbara horisontellt
\end{itemize}

% ═══════════════════════════════════════════════════════════════════════
\section{Sammanfattning}
% ═══════════════════════════════════════════════════════════════════════

AdminDashboard är en central komponent för Roaring Analytics AB att:

\begin{enumerate}
  \item \textbf{Övervaka} alla användare och deras onboarding-status
  \item \textbf{Analysera} fraud detection-resultat för rapportering till Finanspolisen
  \item \textbf{Hantera} support-ärenden från revisionsbyråer
  \item \textbf{Fakturera} kunder baserat på antal användare
  \item \textbf{Kommunicera} via email (enskilt och massutskick)
\end{enumerate}

\subsection{Nästa Steg}
\begin{enumerate}
  \item Testa AdminDashboard i browser (\texttt{http://localhost:5173/admin})
  \item Implementera backend API-endpoints i FastAPI
  \item Integrera SendGrid för produktion
  \item Köp Themesberg-template (25:e november) och migrera till professionell layout
  \item Skapa separerad AdminSidebar och AdminHeader
\end{enumerate}

\vspace{1cm}
\noindent\textbf{Version:} 1.0\\
\textbf{Datum:} \today\\
\textbf{Författare:} Lasse Karagiannis \& Claude AI\\
\textbf{Status:} Mock-version klar, produktion pending

\end{document}
